\part{Geometry, operations, and representations}
\chapter{Extensions of geometric objects}
A sheaf retains more information than its support or its cohomology dimension.
Let $R=\mathbb C[x]$ and $k=R/(x)$.
The sequence
\[
 0\longrightarrow k\xrightarrow{1\mapsto x}R/(x^2)
 \longrightarrow k\longrightarrow0
\]
is a nonsplit extension.
Indeed, $x$ acts nontrivially on its middle term and acts trivially on
$k\oplus k$.
The free resolution $[R\xrightarrow{x}R]$ of $k$ gives
$\operatorname{Ext}^1_R(k,k)=k$.
Thus one extension parameter distinguishes objects with the same two
composition factors.

For a moduli stack $\mathcal M_\gamma$ of objects of charge $\gamma$,
extensions form a correspondence
\[
 \mathcal M_\alpha\times\mathcal M_\beta
 \xleftarrow{p}\mathcal E_{\alpha,\beta}
 \xrightarrow{q}\mathcal M_{\alpha+\beta}.
\]
A coefficient theory with admissible pullback, pushforward, and base
change gives the Hall product $q_*p^*$.
Critical coefficients additionally require orientation isomorphisms
compatible on triple flags.
A framed moduli problem gives a representation when its action
correspondences satisfy the corresponding flag identities.
These geometric operations precede a trace on the representation.

A Calabi--Yau structure supplies a nondegenerate cyclic duality class.
At an object $E$, the tangent complex of its moduli problem is
$R\operatorname{Hom}(E,E)[1]$.
Composition and trace pair tangent vectors with shift $2-d$ for a
$d$-Calabi--Yau category.
Closedness uses the negative-cyclic lift.
Under the hypotheses of shifted-symplectic transgression, this defines
the moduli form.
The ordinary Hochschild cochains have their $E_2$ operations;
an action of another operad or a holomorphic field realization
requires its specified chain-level construction.

\chapter{Central actions and monoidal transport}
For an associative algebra object $A$ in a specified derived coefficient
category, its Hochschild centre is the derived endomorphism object of
its diagonal bimodule.
For a monoidal category $\mathcal M$, its categorical centre
$Z(\mathcal M)$ consists instead of pairs $(X,\beta)$, where
\[
 \beta_Y:X\otimes Y\xrightarrow{\sim}Y\otimes X
\]
is natural in $Y$, compatible with the unit, and satisfies the
half-braiding hexagon.
In strict notation that identity is
$\beta_{Y\otimes Z}=(\id_Y\otimes\beta_Z)(\beta_Y\otimes\id_Z)$.
Associators are inserted in a general monoidal category.
The centre has braiding $c_{(X,\beta),(Y,\gamma)}=\beta_Y$.

\begin{theorem}[Transport of categorical centres]\label{cyarch:centre-transport}
Let $F:\mathcal M\to\mathcal N$ be a strong monoidal equivalence.
Write $J_{X,Y}:FX\otimes FY\to F(X\otimes Y)$ for its tensor constraint.
Choose a monoidal quasi-inverse $H$ and a monoidal natural isomorphism
$\epsilon:FH\to\id_{\mathcal N}$.
Then $F$ induces a braided equivalence
$Z(F):Z(\mathcal M)\to Z(\mathcal N)$.
For $(X,\beta)$ its half-braiding at $Y\in\mathcal N$ is
\[
\begin{split}
 FX\otimes Y
 &\xrightarrow{\id\otimes\epsilon_Y^{-1}}FX\otimes FHY
 \xrightarrow{J_{X,HY}}F(X\otimes HY)\\
 &\xrightarrow{F\beta_{HY}}F(HY\otimes X)
 \xrightarrow{J_{HY,X}^{-1}}FHY\otimes FX
 \xrightarrow{\epsilon_Y\otimes\id}Y\otimes FX.
\end{split}
\]
\end{theorem}
\begin{proof}
Naturality follows by applying naturality of $J$, $\epsilon$, and $\beta$
to each arrow.
For $Y\otimes Z$, monoidality of $\epsilon$ replaces its inverse by the
two inverses and the tensor constraint of $FH$.
The pentagon for $J$ then identifies the transported hexagon with
$F$ applied to the hexagon for $\beta$.
The same argument with the unit constraints proves the unit identity.
The constraints $J$ are central morphisms by this formula.
For two central objects their braiding is the half-braiding evaluated
on the second object, so the construction preserves braiding.
Applying the same construction to $H$ gives an inverse up to the
monoidal unit and counit isomorphisms.
\end{proof}

A Morita equivalence of associative algebras supplies an equivalence of
module categories.
To use this theorem, those categories must have specified tensor
products and the equivalence must have coherent tensor and unit
constraints.
A realization of the Hochschild centre by a categorical centre is a
further comparison with its own source and target.

\chapter{Field actions and amplitudes}
Fix a spacetime, fields, action, gauge symmetry, boundary condition,
coefficient category, and approximation regime.
A quantization must satisfy the quantum master equation on its stated
domain before it supplies quantum observables.
A coherent action on the boundary has target the independently defined
boundary centre:
\[
 \beta:\operatorname{Obs}_{\mathrm{bulk}}\big|_\partial
 \longrightarrow Z_{\mathrm{ord}}(A_\partial).
\]
An equivalence of this map on a specified sector is a bulk reconstruction
statement.
A supplied formal factorization envelope is a different object until
its realization map to the spacetime observables is constructed.

Hall convolution produces a positive algebra.
A compatible coproduct gives tensor products of representations.
A paired negative half and Cartan part permit a completed double when
the triangular decomposition and canonical tensors converge.
Exchange requires the corresponding intertwiners and their hexagon
identities.
Theorem~\ref{cyarch:centre-transport} then applies to actual monoidal
comparisons of the representation categories.

A trace closes an action into an amplitude only on its declared domain.
Its cyclicity, grading, value line, and normalization are additional data.
A Borcherds character computes signed graded dimensions of its root
algebra.
Its identification with a geometric Hall amplitude requires root states,
orientation maps, and a trace comparison on those states.
